\section{Determinant specification of \texorpdfstring{\(Q_D\)}{Q-D}}
\label{app:quartic-specification}

This appendix fixes every convention needed to recover the quartic without
printing its expanded degree-\(36\) coefficient table.

\subsection{Gauge and monomial order}

List all exponent vectors
\[
 (e_0,e_1,e_2,e_3),\qquad e_i\ge0,\qquad
 e_0+e_1+e_2+e_3=4,
\]
in descending lexicographic order.  Delete
\[
 (4,0,0,0),\quad(3,1,0,0),\quad(3,0,1,0),\quad(3,0,0,1).
\]
The validity and uniqueness of this deletion are proved by the gauge block
\cref{eq:gauge-block,eq:gauge-determinant}.  The remaining ordered monomials
are
\begin{align*}
m_0&=u_0^2u_1^2,&
m_1&=u_0^2u_1u_2,&
m_2&=u_0^2u_1u_3,\\
m_3&=u_0^2u_2^2,&
m_4&=u_0^2u_2u_3,&
m_5&=u_0^2u_3^2,\\
m_6&=u_0u_1^3,&
m_7&=u_0u_1^2u_2,&
m_8&=u_0u_1^2u_3,\\
m_9&=u_0u_1u_2^2,&
m_{10}&=u_0u_1u_2u_3,&
m_{11}&=u_0u_1u_3^2,\\
m_{12}&=u_0u_2^3,&
m_{13}&=u_0u_2^2u_3,&
m_{14}&=u_0u_2u_3^2,\\
m_{15}&=u_0u_3^3,&
m_{16}&=u_1^4,&
m_{17}&=u_1^3u_2,\\
m_{18}&=u_1^3u_3,&
m_{19}&=u_1^2u_2^2,&
m_{20}&=u_1^2u_2u_3,\\
m_{21}&=u_1^2u_3^2,&
m_{22}&=u_1u_2^3,&
m_{23}&=u_1u_2^2u_3,\\
m_{24}&=u_1u_2u_3^2,&
m_{25}&=u_1u_3^3,&
m_{26}&=u_2^4,\\
m_{27}&=u_2^3u_3,&
m_{28}&=u_2^2u_3^2,&
m_{29}&=u_2u_3^3,\\
m_{30}&=u_3^4.&&&
\end{align*}
The normalization is \(q_0=1\); equivalently,
\([u_0^2u_1^2]Q_D=1\).

\subsection{The \texorpdfstring{\(60\times31\)}{60 by 31} matrix}

For a carrier root \(d\), write the corresponding line as
\[
 (s,t,a(d)s+c(d)t,b(d)s+dt).
\]
For every \(m_j\), expand
\begin{equation}
 m_j(s,t,a(d)s+c(d)t,b(d)s+dt)
 =\sum_{\tau=0}^4r_{\tau j}(d)s^{4-\tau}t^\tau.
\label{eq:binary-restriction}
\end{equation}
Reduce \(r_{\tau j}(d)\) modulo the monic carrier \(\Acarrier(d)\):
\[
 r_{\tau j}(d)\equiv
 \sum_{\rho=0}^{11}M_{12\tau+\rho,j}d^\rho
 \pmod{\Acarrier(d)}.
\]
This fixes the row convention
\begin{equation}
 \operatorname{row}=12(\text{binary degree})+(\text{carrier degree}),
\label{eq:row-convention}
\end{equation}
with both degrees increasing.  The matrix
\[
 M_D=(M_{rj})_{\substack{0\le r<60\\0\le j<31}}
\]
is therefore determined by the surface, chart shape, carrier, gauge, and
orders above.

\subsection{Pivot and Cramer's rule}

Let
\[
 \cR=\{0,1,\ldots,10,12,\ldots,20,24,\ldots,29,36,37,38,48\}.
\]
Define
\[
 N_D=(M_{rj})_{\substack{r\in\cR\\1\le j\le30}},
 \qquad
 b_D=(M_{r0})_{r\in\cR}.
\]
The exact pivot certificate proves \(\det N_D\ne0\).  If
\(N_D^{(j)}(b_D)\) denotes \(N_D\) with its \(j\)-th column replaced by
\(b_D\), then
\begin{equation}
 q_j=-\frac{\det N_D^{(j)}(b_D)}{\det N_D},
 \qquad 1\le j\le30.
\label{eq:cramer-Q}
\end{equation}
\Cref{eq:QD-definition,eq:cramer-Q}, together with the monomial
list, define \(Q_D\) exactly.

The pivot determinant has nonzero norm \(3\) modulo \(7\).  The exact orbit
replay checks all \(36\) conjugate determinant values and all
\(36\cdot60=2160\) restriction equations.  Stable internal fingerprints are
\begin{center}
\begin{tabular}{@{}l@{}}
pivot determinant coefficient digest:\\
\code{1af5bfdc9b2f945094835fd81281305fb84dfc8208cb542874f2803420cd3a9e}\\[3pt]
canonical solution digest:\\
\code{eb9e803e16f5623647843dd7636fe3d511af60f8f31c453ea6826cd3e4d25573}
\end{tabular}
\end{center}
The fingerprints audit the exact arrays.  They are not mathematical
substitutes for gauge invertibility, the Hilbert--90 upper bound, pivot
nonvanishing, or the complete restriction replay.
